\documentclass[12pt,a4paper]{article}

% ------------------------------------------------
% Codificación y lenguaje
% ------------------------------------------------
\usepackage[utf8]{inputenc}      % Codificación UTF-8
\usepackage[T1]{fontenc}
\usepackage[spanish]{babel}
\usepackage{lmodern}

% ------------------------------------------------
% Márgenes y formato
% ------------------------------------------------
\usepackage[a4paper,margin=2.5cm]{geometry}
\usepackage{setspace}
\onehalfspacing
% ------------------------------------------------
% Matemáticas
% ------------------------------------------------
\usepackage{amsmath,amssymb,amsthm,mathtools}
\usepackage{bm}

% Entornos de teorema
\newtheorem{teorema}{Teorema}[section]
\newtheorem{proposicion}[teorema]{Proposición}
\newtheorem{lema}[teorema]{Lema}
\newtheorem{corolario}[teorema]{Corolario}
\theoremstyle{definition}
\newtheorem{definicion}[teorema]{Definición}
\theoremstyle{remark}
\newtheorem{observacion}[teorema]{Observación}

% ------------------------------------------------
% Gráficos y tablas
% ------------------------------------------------
\usepackage{graphicx}
\usepackage{float}
\usepackage{booktabs}
\usepackage{array}
\usepackage{multirow}

% ------------------------------------------------
% Código MATLAB
% ------------------------------------------------
\usepackage{listings}
\usepackage{xcolor}

\definecolor{codegray}{gray}{0.95}
\definecolor{codegreen}{rgb}{0,0.5,0}
\definecolor{codeblue}{rgb}{0,0,0.6}

\lstdefinestyle{matlabstyle}{
    language=Matlab,
    backgroundcolor=\color{codegray},
    commentstyle=\color{codegreen},
    keywordstyle=\color{codeblue}\bfseries,
    numberstyle=\tiny\color{gray},
    stringstyle=\color{red},
    basicstyle=\ttfamily\small,
    breaklines=true,
    frame=single,
    numbers=left,
    numbersep=5pt,
    tabsize=2,
    showstringspaces=false
}

\lstset{style=matlabstyle}

% ------------------------------------------------
% Hipervínculos
% ------------------------------------------------
\usepackage[hidelinks]{hyperref}

% ------------------------------------------------
% Comandos útiles
% ------------------------------------------------
\newcommand{\R}{\mathbb{R}}
\newcommand{\N}{\mathbb{N}}
\newcommand{\Ogrande}{\mathcal{O}}
\newcommand{\eps}{\varepsilon}

% ------------------------------------------------
% Documento
% ------------------------------------------------
\title{Derivación Numérica}
\author{Nico Iko 777}
\date{\today}

\begin{document}

\maketitle
\tableofcontents
\newpage

Cuidado: La guía no está exenta de errores !!!


\section{Derivación numérica}

En muchos problemas científicos la función $f(x)$ no está disponible de forma analítica,
sino únicamente a través de valores en una malla

\[
x_0, x_1, x_2, \dots , x_n .
\]

En estas situaciones la derivada debe aproximarse utilizando los valores conocidos
de la función. Este procedimiento se conoce como \textbf{derivación numérica}.

Sea $h$ el tamaño de paso de la malla uniforme

\[
x_i = x_0 + ih .
\]

La idea fundamental consiste en aproximar la derivada mediante combinaciones
lineales de valores de la función evaluados en puntos cercanos.


En cálculo numérico:
\begin{itemize}
\item No podemos tomar el límite.
\item Solo disponemos de valores discretos $f(x_i)$.
\item Los datos pueden provenir de mediciones experimentales.
\end{itemize}

Por tanto, sustituimos el límite por un cociente incremental con $h$ pequeño pero finito.\\

Muchos métodos numéricos dependen de un parámetro $h$ (tamaño de paso). 
Si conocemos cómo se comporta el error respecto a $h$, podemos combinar 
aproximaciones con distintos valores de $h$ para eliminar el término 
dominante del error.

\subsection{Aproximación básica}

A partir de la definición de derivada

\[
f'(x)=\lim_{h\to0}\frac{f(x+h)-f(x)}{h}
\]

se obtiene la aproximación más simple

\[
f'(x) \approx \frac{f(x+h)-f(x)}{h}.
\]

Esta aproximación tiene error de truncamiento

\[
O(h).
\]

El objetivo de los métodos de derivación numérica es construir fórmulas con
errores más pequeños utilizando más puntos de la malla.


%%%%%%%%%%%%%%%%%%%%%

\section{Fórmulas de diferencias para la primera derivada}

Las fórmulas de derivación numérica se obtienen a partir de expansiones de Taylor
de la función alrededor del punto de interés.


Sea $f$ suficientemente suave. Entonces

\[
f(x+h) = f(x) + h f'(x) + \frac{h^2}{2}f''(x) + \frac{h^3}{6}f^{(3)}(x) + \cdots
\]

A partir de esta expansión se diseñan las fórmulas.

\subsection{Diferencia hacia adelante}

Si se conocen los valores de la función en $x$ y $x+h$, se obtiene

\[
f'(x) \approx
\frac{f(x+h)-f(x)}{h},
\]

con error

\[
O(h).
\]

Una aproximación de mayor precisión utiliza tres puntos

\[
x,\; x+h,\; x+2h
\]

y está dada por

\[
f'(x) \approx
\frac{-3f(x)+4f(x+h)-f(x+2h)}{2h},
\]

con error

\[
O(h^2).
\]

\subsection{Diferencia hacia atrás}

De forma análoga, utilizando los puntos

\[
x,\; x-h,\; x-2h
\]

se obtiene

\[
f'(x) \approx
\frac{3f(x)-4f(x-h)+f(x-2h)}{2h},
\]

con error

\[
O(h)= \frac{h^2 f^{(3)}(\xi)}{2}
\]




\subsection{Diferencia centrada}

Las fórmulas centradas utilizan información a ambos lados del punto.

La aproximación más común es

\[
f'(x) \approx
\frac{f(x+h)-f(x-h)}{2h},
\]

con error


\[
O(h^2)= \frac{h^2 f^{(3)}(\xi)}{6}
\]




\subsubsection*{Ejemplo 2: Comparación}

Sea $f(x)=e^x$, en $x=0$.

\[
f'(0)=1.
\]

\textbf{Diferencia hacia adelante:}

\[
D_h^+ f(0)=\frac{e^h-1}{h}.
\]

Usando Taylor:

\[
\frac{e^h-1}{h}
=
1 + \frac{h}{2} + O(h^2).
\]

Error $\sim \frac{h}{2}$.

\bigskip

\textbf{Diferencia centrada:}

\[
D_h^c f(0)=\frac{e^h-e^{-h}}{2h}.
\]

Desarrollando:

\[
=
1 + \frac{h^2}{6} + O(h^4).
\]

Error $\sim \frac{h^2}{6}$.

\bigskip

Para $h=10^{-2}$:

\[
\text{Error adelante} \sim 5\times 10^{-3},
\]
\[
\text{Error centrado} \sim 1.6\times 10^{-5}.
\]

La mejora es significativa.

\subsection{Fórmula centrada de orden $O(h^4)$}

Si se utilizan cinco puntos

\[
x-2h,\;x-h,\;x,\;x+h,\;x+2h,
\]

se obtiene la aproximación

\[
f'(x) \approx
\frac{-f(x+2h)+8f(x+h)-8f(x-h)+f(x-2h)}{12h},
\]

con error de truncamiento


\[
O(h^4)= \frac{h^4 f^{(5)}(\xi)}{30}
\]

Esta fórmula se obtiene combinando las expansiones de Taylor de
$f(x\pm h)$ y $f(x\pm2h)$ de forma que se cancelen los términos
de error de menor orden.





\subsection*{Diferencia hacia adelante de orden $O(h^4)$}

Cuando no se dispone de puntos a la izquierda del punto de evaluación,
se pueden utilizar fórmulas hacia adelante. Usando cinco puntos se obtiene
una aproximación de orden cuatro.

Sea $h>0$ y supongamos que conocemos los valores de la función en los puntos

\[
x,\; x+h,\; x+2h,\; x+3h,\; x+4h.
\]

La derivada puede aproximarse mediante

\[
f'(x) \approx 
\frac{-25f(x) + 48f(x+h) - 36f(x+2h) + 16f(x+3h) - 3f(x+4h)}{12h}.
\]



\subsubsection*{Ejemplo}

Sea

\[
f(x)=\sin x
\]

y queremos aproximar la derivada en $x=1$.

Sabemos que

\[
f'(x)=\cos x.
\]

Tomando $h=0.1$ y aplicando la fórmula anterior obtenemos

\[
f'(1) \approx
\frac{-25f(1)+48f(1.1)-36f(1.2)+16f(1.3)-3f(1.4)}{12(0.1)}.
\]

El valor exacto es

\[
f'(1)=\cos(1).
\]

La aproximación numérica presenta un error del orden $O(h^4)$,
mucho menor que las fórmulas de orden $O(h)$ o $O(h^2)$.

\subsection{Número de puntos y orden de precisión}

\begin{center}
\begin{tabular}{ccc}
\hline
Método & Puntos utilizados & Error \\ \hline
Forward & $x,x+h$ & $O(h)$ \\
Forward & $x,x+h,x+2h$ & $O(h^2)$ \\
Backward & $x,x-h,x-2h$ & $O(h^2)$ \\
Centered & $x-h,x,x+h$ & $O(h^2)$ \\
Centered & $x-2h,x-h,x,x+h,x+2h$ & $O(h^4)$ \\
Centered & $x-3h,\dots,x+3h$ & $O(h^6)$ \\
\hline
\end{tabular}
\end{center}


\section{Extrapolación de Richardson}

La extrapolación de Richardson es una técnica que permite aumentar el orden
de precisión de una aproximación numérica utilizando cálculos con diferentes
valores del paso $h$.

Supongamos que una aproximación de la derivada tiene la forma

\[
D(h) = f'(x) + C h^p + O(h^{p+1}).
\]

Si se calcula la misma aproximación con paso $h/2$

\[
D(h/2) = f'(x) + C\left(\frac{h}{2}\right)^p + O(h^{p+1}),
\]

entonces es posible eliminar el término dominante del error.

La extrapolación de Richardson está dada por

\[
D_R =
\frac{2^p D(h/2)-D(h)}{2^p-1}.
\]

El resultado es una aproximación con error

\[
O(h^{p+1}).
\]

Esta técnica se utiliza frecuentemente para mejorar la precisión de métodos
de diferencias finitas sin aumentar el número de puntos de la malla.


\subsection{Ejemplo de extrapolación de Richardson}

Consideremos la derivada de

\[
f(x)=\sin x
\]

en el punto $x=0$. Sabemos que

\[
f'(0)=\cos(0)=1.
\]

Usaremos la fórmula centrada

\[
D_h f(x)=\frac{f(x+h)-f(x-h)}{2h},
\]

la cual tiene error de orden $O(h^2)$.

\subsubsection*{Paso 1: aproximación con $h=0.1$}

\[
D_{0.1}f(0)
=
\frac{\sin(0.1)-\sin(-0.1)}{0.2}
\approx 0.998334.
\]

\subsubsection*{Paso 2: aproximación con $h/2=0.05$}

\[
D_{0.05}f(0)
=
\frac{\sin(0.05)-\sin(-0.05)}{0.1}
\approx 0.999583.
\]

\subsubsection*{Paso 3: extrapolación de Richardson}

Como el método es de orden $p=2$, aplicamos

\[
D_R
=
\frac{2^p D(h/2)-D(h)}{2^p-1}.
\]

Entonces

\[
D_R
=
\frac{4(0.999583)-0.998334}{3}
\approx 0.999999.
\]

\subsubsection*{Conclusión}

La aproximación original tenía error de orden $O(h^2)$, mientras que
la extrapolación de Richardson produce una aproximación de orden $O(h^4)$,
mucho más precisa.


\section{Ruido, estabilidad y precisión}

\subsection{Ruido}

\begin{definicion}[Ruido]
Supongamos que en lugar de conocer los valores exactos $f(x)$ disponemos de datos
\[
\tilde f(x)=f(x)+\eta(x),
\]
donde $\eta(x)$ es una perturbación pequeña. Esta perturbación se denomina
\textbf{ruido}.
\end{definicion}

El ruido puede deberse a errores de medición, redondeo numérico o
modelización incompleta.

La derivación numérica tiende a amplificar el ruido. Por ejemplo,

\[
\frac{\tilde f(x+h)-\tilde f(x)}{h}
=
\frac{f(x+h)-f(x)}{h}
+
\frac{\eta(x+h)-\eta(x)}{h}.
\]

Si el ruido es del orden $\delta$, entonces el término perturbador es
aproximadamente

\[
\frac{\delta}{h}.
\]

Por tanto, si $h$ es muy pequeño el ruido puede amplificarse significativamente.

\subsection{Precisión}

La \textbf{precisión} de una fórmula de derivación se refiere al orden de su error
de truncamiento. Si

\[
f'(x)=D_h f(x)+Ch^p+O(h^{p+1}),
\]

decimos que el método es de orden $p$.

Reducir $h$ a la mitad disminuye el error aproximadamente por un factor $2^p$.

Ejemplos típicos:

\begin{itemize}
\item Diferencia hacia adelante: orden $1$.
\item Diferencia centrada: orden $2$.
\item Fórmulas centradas de cinco puntos: orden $4$.
\end{itemize}

\subsection{Estabilidad}

Un método numérico es \textbf{estable} si pequeñas perturbaciones en los datos
producen pequeñas variaciones en el resultado. En derivación numérica esto
significa que

\[
\|\tilde f-f\|\ \text{pequeño}
\quad \Rightarrow \quad
\|D_h\tilde f-D_h f\|\ \text{pequeño}.
\]

Las derivadas amplifican las componentes de alta frecuencia. Por ejemplo,
si $f(x)=\sin(\omega x)$ entonces

\[
f'(x)=\omega\cos(\omega x).
\]

Las frecuencias altas ($\omega$ grande) se amplifican por un factor $\omega$,
lo cual explica por qué el ruido puede afectar seriamente a la derivación
numérica.


\subsubsection*{Ejemplo 3: Inestabilidad con datos oscilatorios}

Sea

\[
f(x)=x^2 + 10^{-3}\sin(1000x).
\]

La parte oscilatoria tiene pequeña amplitud, pero:

\[
\frac{d}{dx}\big(10^{-3}\sin(1000x)\big)
=
10^{-3}\cdot 1000\cos(1000x)
=
1\cos(1000x).
\]

Una perturbación pequeña en amplitud produce una contribución grande en la derivada.

\subsection{Relación entre ruido, precisión y estabilidad}

El error total de una fórmula de derivación puede modelarse como

\[
E(h)\approx C_1 h^p + \frac{C_2}{h}\,\delta,
\]

donde $h^p$ representa el error de truncamiento y $\delta/h$ la amplificación
del ruido.

En consecuencia, disminuir $h$ reduce el error de truncamiento pero aumenta
la influencia del ruido. En la práctica existe un valor intermedio de $h$
que minimiza el error total.

\begin{itemize}
\item La \textbf{precisión} describe cómo mejora el método cuando $h\to0$.
\item La \textbf{estabilidad} mide la sensibilidad frente a perturbaciones.
\item El \textbf{ruido} es inevitable en datos reales.
\end{itemize}


\section{Mas ejemplos}


\subsection{Ejemplo}

Queremos aproximar

\[
f'(0), \quad f(x)=e^x.
\]

Usamos diferencia centrada

\[
D_h = \frac{f(h)-f(-h)}{2h}.
\]

Con $h=0.1$

\[
D_{0.1} \approx 1.00167
\]

Con $h=0.05$

\[
D_{0.05} \approx 1.00042
\]

Como el método es de orden $p=2$:

\[
D_R = \frac{4D_{0.05}-D_{0.1}}{3}
\]

\[
D_R \approx 1.00000
\]

que coincide prácticamente con el valor exacto

\[
f'(0)=1.
\]

\newpage

\section{Diferencias finitas}

Las fórmulas de derivación numérica pertenecen a la familia de los métodos
de \textbf{diferencias finitas}. En estos métodos las derivadas se aproximan
mediante combinaciones lineales de valores de la función en una malla discreta.

\subsection{Malla}

Se considera una malla uniforme

\[
x_i = x_0 + ih, \qquad i=0,1,\dots,n,
\]

donde $h$ es el tamaño de paso.

\subsection{Stencil}

El conjunto de puntos utilizados para construir una aproximación se denomina
\textbf{stencil}. Por ejemplo,

\[
x_{i-1},x_i,x_{i+1}
\]

define el stencil de la fórmula centrada de orden $O(h^2)$.

\subsection{Diferencias finitas y diferencias centradas}

Las \textbf{diferencias finitas} constituyen el marco general para aproximar
derivadas en una malla discreta.

Dentro de este marco existen diferentes tipos de fórmulas:

\begin{itemize}

\item \textbf{Diferencias hacia adelante}, que utilizan puntos a la derecha
del punto de evaluación.

\item \textbf{Diferencias hacia atrás}, que utilizan puntos a la izquierda.

\item \textbf{Diferencias centradas}, que utilizan puntos a ambos lados.

\end{itemize}

Las fórmulas centradas suelen ser más precisas debido a la simetría del stencil,
la cual permite cancelar términos adicionales en la expansión de Taylor.

\subsection{Ejemplo paso a paso de diferencias finitas}

Consideremos la función

\[
f(x)=x^2.
\]

Sabemos que su derivada exacta es

\[
f'(x)=2x.
\]

Nuestro objetivo es aproximar la derivada utilizando diferencias finitas.

\subsubsection*{Paso 1: Construir una malla}

Tomemos el intervalo

\[
[0,1]
\]

y dividámoslo en puntos igualmente espaciados con paso

\[
h=0.25.
\]

La malla es

\[
x_0=0,\quad x_1=0.25,\quad x_2=0.5,\quad x_3=0.75,\quad x_4=1.
\]

\subsubsection*{Paso 2: Evaluar la función}

Calculamos los valores de la función:

\[
\begin{array}{c|c}
x & f(x)=x^2 \\
\hline
0 & 0 \\
0.25 & 0.0625 \\
0.5 & 0.25 \\
0.75 & 0.5625 \\
1 & 1
\end{array}
\]

\subsubsection*{Paso 3: Aplicar diferencias centradas}

Usamos la fórmula

\[
f'(x_i)\approx
\frac{f(x_{i+1})-f(x_{i-1})}{2h}.
\]

\subsubsection*{Paso 4: Calcular derivadas aproximadas}

Para $x_1=0.25$

\[
f'(0.25)\approx
\frac{f(0.5)-f(0)}{2h}
=
\frac{0.25-0}{0.5}
=
0.5.
\]

Para $x_2=0.5$

\[
f'(0.5)\approx
\frac{f(0.75)-f(0.25)}{0.5}
=
\frac{0.5625-0.0625}{0.5}
=
1.
\]

Para $x_3=0.75$

\[
f'(0.75)\approx
\frac{f(1)-f(0.5)}{0.5}
=
\frac{1-0.25}{0.5}
=
1.5.
\]

\subsubsection*{Paso 5: Comparación}

\[
\begin{array}{c|c|c}
x & \text{Derivada numérica} & \text{Derivada exacta} \\
\hline
0.25 & 0.5 & 0.5 \\
0.5 & 1 & 1 \\
0.75 & 1.5 & 1.5
\end{array}
\]

En este caso el resultado coincide exactamente porque $f(x)=x^2$
es un polinomio de grado bajo y la fórmula centrada reproduce
exactamente su derivada.

\newpage
\subsection{Errores en la derivación numérica}

Las aproximaciones numéricas de derivadas presentan dos tipos principales de error:
el \textbf{error de truncamiento} y el \textbf{error de redondeo}.

Si una aproximación satisface

\[
E(h) = C h^p + O(h^{p+1}),
\]

se dice que el método es de \textbf{orden $p$}.

\subsubsection*{Error de truncamiento}

Este error aparece al truncar la serie de Taylor utilizada para derivar
las fórmulas de diferencias.

Su comportamiento típico es

\[
E_T \sim C h^p.
\]

Por tanto, disminuir el tamaño de paso $h$ reduce el error de truncamiento.

\subsubsection*{Error de redondeo}

En aritmética de precisión finita aparecen errores debido a la representación
numérica de los números. Para las fórmulas de derivación suele estimarse como

\[
E_R \sim \frac{\varepsilon_{\text{mach}}}{h},
\]

donde $\varepsilon_{\text{mach}}$ es la precisión de máquina.

\subsubsection*{Error total}

El error total puede aproximarse por

\[
E(h) \approx C_1 h^p + \frac{C_2 \varepsilon_{\text{mach}}}{h}.
\]

Si $h$ es demasiado grande domina el error de truncamiento, mientras que si
$h$ es demasiado pequeño domina el error de redondeo. En consecuencia, existe
un valor intermedio de $h$ que minimiza el error total.

\subsection{Limitaciones de la derivación numérica}

La derivación numérica puede producir resultados poco fiables en
las siguientes situaciones:

\begin{itemize}

\item Cuando la función no es diferenciable.

\item Cuando los datos contienen ruido experimental.

\item Cuando los puntos de la malla están demasiado separados.

\item En presencia de discontinuidades o cambios bruscos.

\end{itemize}

\subsection{Aplicaciones}

Las derivadas numéricas aparecen en numerosos problemas científicos y de ingeniería:

\begin{itemize}

\item Métodos para ecuaciones diferenciales ordinarias (Euler, Runge--Kutta).

\item Métodos para ecuaciones en derivadas parciales (métodos de diferencias finitas).

\item Optimización numérica (cálculo de gradientes).

\item Procesamiento de señales e imágenes.

\item Mecánica computacional y dinámica de fluidos.

\item Transferencia de calor y difusión.

\end{itemize}



\section{Apendice}
\subsection{Diferencias hacia adelante de alto orden}

Para aproximar derivadas cuando sólo se conocen valores de la función en una malla uniforme
\[
x, \; x+h, \; x+2h, \; x+3h, \; x+4h, \dots
\]
pueden construirse fórmulas de diferencias finitas utilizando expansiones de Taylor.
La idea consiste en combinar valores de la función de forma que se cancelen los términos
de error de menor orden.

\subsection{Expansiones de Taylor}

Supongamos que $f$ es suficientemente suave. Las expansiones alrededor de $x$ son

\begin{align}
f(x+h) &= f(x) + hf'(x) + \frac{h^2}{2}f''(x) + \frac{h^3}{6}f'''(x)
+ \frac{h^4}{24}f^{(4)}(x) + O(h^5),\\
f(x+2h) &= f(x) + 2hf'(x) + \frac{(2h)^2}{2}f''(x) + \frac{(2h)^3}{6}f'''(x)
+ \frac{(2h)^4}{24}f^{(4)}(x) + O(h^5),\\
f(x+3h) &= f(x) + 3hf'(x) + \frac{(3h)^2}{2}f''(x) + \frac{(3h)^3}{6}f'''(x)
+ \frac{(3h)^4}{24}f^{(4)}(x) + O(h^5),\\
f(x+4h) &= f(x) + 4hf'(x) + \frac{(4h)^2}{2}f''(x) + \frac{(4h)^3}{6}f'''(x)
+ \frac{(4h)^4}{24}f^{(4)}(x) + O(h^5).
\end{align}

\subsection{Combinación lineal}

Buscamos una combinación lineal

\[
a_0 f(x) + a_1 f(x+h) + a_2 f(x+2h) + a_3 f(x+3h) + a_4 f(x+4h)
\]

que aproxime $f'(x)$. Sustituyendo las expansiones de Taylor obtenemos

\begin{align}
&a_0 f(x) + a_1 f(x+h) + a_2 f(x+2h) + a_3 f(x+3h) + a_4 f(x+4h) \\
&= (a_0+a_1+a_2+a_3+a_4)f(x) \\
&+ h(a_1+2a_2+3a_3+4a_4)f'(x) \\
&+ \frac{h^2}{2}(a_1+4a_2+9a_3+16a_4)f''(x) \\
&+ \frac{h^3}{6}(a_1+8a_2+27a_3+64a_4)f'''(x) \\
&+ \frac{h^4}{24}(a_1+16a_2+81a_3+256a_4)f^{(4)}(x) + O(h^5).
\end{align}

Para obtener una aproximación de orden $O(h^4)$ se imponen las condiciones

\begin{align}
a_0+a_1+a_2+a_3+a_4 &= 0, \\
a_1+2a_2+3a_3+4a_4 &= 1, \\
a_1+4a_2+9a_3+16a_4 &= 0, \\
a_1+8a_2+27a_3+64a_4 &= 0, \\
a_1+16a_2+81a_3+256a_4 &= 0.
\end{align}

Resolviendo este sistema se obtiene

\[
a_0=-\frac{25}{12h},\quad
a_1=\frac{48}{12h},\quad
a_2=-\frac{36}{12h},\quad
a_3=\frac{16}{12h},\quad
a_4=-\frac{3}{12h}.
\]

Por tanto, la aproximación final es

\[
f'(x) \approx 
\frac{-25f(x)+48f(x+h)-36f(x+2h)+16f(x+3h)-3f(x+4h)}{12h},
\]

con error

\[
O(h^4).
\]


\end{document}